\section{弱楕円gluingのcoherent反標準関手と収縮圏}\label{sec:coherent-contraction}
この節では第\ref{sec:zero-walls}節の標準弱楕円gluingを固定する．
$S$は滑らかな射影有理弱del Pezzo曲面，$E\in|-K_S|$は滑らかな楕円曲線，
$d=K_S^2>0$，$q\in\Pic^0(E)$である．$q$の位数に制限を課さない．
$e\in\cD_q$は構造exceptional対象，$s=s_q$は反Serre作用とし，
$F_qs\simeq\Psi_qF_q$と書く．本節の結論は任意の抽象的な非可換曲面ではなく，
この指定gluingとその自然な境界triangleについてのものである．

\subsection{境界の正次数性からの有限生成}
\begin{lemma}[次数付きNakayamaと境界切断]\label{lem:graded-boundary-finiteness}
次の二つが成り立つ．
\begin{enumerate}
\item 非負次数付き代数$R$の正次数中心元$g$と，下に有界な次数付き右加群$N$について，
$N/gN$が有限生成なら$N$も有限生成である．この主張は$N$上の$g$の単射性を要求しない．
\item $\Psi$を楕円曲線上の平行移動と正次数のline bundleによるtensorの合成とし，
$B=\bigoplus_{n\ge0}H^0(E,\Psi^n\OO_E)$とする．任意の$V\in\Db(\coh E)$について
\[
 C^j(V)=\bigoplus_{n\ge0}H^j(E,\Psi^nV)
\]
は右$B$有限生成で，非零となる$j$は有限個である．
\end{enumerate}
\end{lemma}
\begin{proof}
(1)では$N/gN$の有限個の同次生成元を$N$へ持ち上げ，それらが生成する部分加群を$L$とする．
$N=L+gN$なので，$N$の最低次数からの帰納法により$N=L$となる．

(2)をまずcoherent sheaf $V$について示す．十分大きい$a$を選び，
$\OO_E^{\oplus h}\twoheadrightarrow\Psi^aV$の核を$K$とする．
平行移動は次数とHarder--Narasimhan slopesを変えず，$\Psi^m$は全slopesを
$m\deg(\Psi\OO_E)$だけ増加させる．従って$n\gg0$で
$H^1(E,\Psi^nK)=0$であり，$\Psi^n$でtwistした評価列から
\[
 B_n^{\oplus h}\twoheadrightarrow H^0(E,\Psi^{a+n}V)
\]
を得る．この写像は次数$a$の切断への右$B_n$作用である．
有限個の低次数成分を加えれば$C^0(V)$の有限生成性が従い，$C^1(V)$は有限次元である．
$B$の左右noetherian性はtwisted homogeneous coordinate ringの既知定理による．
平行移動は数値的divisor群に恒等的に作用し，正次数line bundleはampleなので，
Keelerの判定と左右noetherian性が適用される
\cite[Theorem 1.2(1),(2)]{KeelerAmple}．
一般の有界複体では，有限個のcohomology sheafに対するhypercohomology spectral sequenceを用いる．
$n\gg0$ではその高次sheaf cohomologyが消え，
$H^j(E,\Psi^nV)=H^0(E,\Psi^nH^jV)$となる．
この同定は$B$作用と両立し，有限個の低次数を補って結論を得る．
曲線のcohomological dimensionが1であることから，$j$の有界性も従う．
\end{proof}

\begin{theorem}[弱gluingにおけるcoherent導来切断]\label{thm:weak-coherent-sections}
標準弱楕円gluingでは，反標準環
\[
 R_q=\bigoplus_{n\ge0}\Hom(e,s^ne)
\]
は局所有限な連結左右noetherian代数で，中心正則元$g\in(R_q)_1$を持ち，
$R_q/gR_q\simeq B(E,\Psi_q)$である．任意の$X\in\cD_q$について
\[
 N^j(X)=\bigoplus_{n\ge0}\Hom(e,s^nX[j])
\]
は右$R_q$有限生成で，非零となる$j$は有限個である．
自然なdg関手
\[
 \Phi_q=\pi\mathbf R\Gamma_*:\cD_q\longrightarrow\Db(\qgr R_q)
\]
が定まり，$H^j\Phi_qX=\pi N^j(X)$，$\Phi_qe=\pi R_q$，
$\Phi_qs\simeq(1)\Phi_q$を満たす．さらに
\begin{equation}\label{eq:weak-coherent-kernel-final}
 \boxed{\ker\Phi_q=\mathcal N_q=\iota_q\ker(Rf_*).}
\end{equation}
従って，標準弱楕円gluingでは命題\ref{prop:weak-coherent-kernel}のcoherent有限性の条件が成立する．
\end{theorem}
\begin{proof}
\emph{反標準環．}
境界triangleを$s^ne$に適用すると
\[
 \RHom(e,s^{n-1}e)\longrightarrow\RHom(e,s^ne)
 \longrightarrow R\Gamma(E,\Psi_q^n\OO_E)\longrightarrow
\]
を得る．$n\ge1$では最後のline bundleの次数は$nd>0$なので，そのcohomologyは次数0のみである．
$e$のexceptionalityを初期値として帰納すると，
$\Ext^j(e,s^ne)=0$ $(j\ne0,n\ge0)$と
\[
 0\longrightarrow(R_q)_{n-1}\xrightarrow{\ g\ }(R_q)_n
 \longrightarrow B(E,\Psi_q)_n\longrightarrow0\quad(n\ge1)
\]
を得る．積との両立性と$g$の中心性は，境界の自然変換によるものであり，
Rains\cite[Proposition 5.1と式(5.3)]{Filtered}と同じ構成である．
次数0も$k$なので商は$B(E,\Psi_q)$となる．

$g$進filtrationの随伴次数付き環は$B(E,\Psi_q)[g]$で左右noetherianである．
同次idealのinitial idealの有限生成元を持ち上げて$g$進次数で還元する．
固定した全次数では$g$進filtrationは有限なので還元は有限回で終わる．
従って同次idealは有限生成であり，一般のidealには全次数filtrationを用いて
同じ結論を得る．これが$R_q$の左右noetherian性を与える．

\emph{次数付きExtの有限生成と有界性．}
一般の$X$に対する境界triangleから，$n\ge1$で
\begin{equation}\label{eq:weak-central-exact}
 N^j(X)_{n-1}\xrightarrow{\ g\ }N^j(X)_n
 \longrightarrow C^j(F_qX)_n\longrightarrow N^{j+1}(X)_{n-1}
\end{equation}
が完全である．これは境界restrictionと積の自然性により次数付き右$R_q$加群の列であり，
$C^j(F_qX)$への作用は$R_q\twoheadrightarrow B(E,\Psi_q)$を経由する．
従って$(N^j(X)/gN^j(X))_{\ge1}$は$C^j(F_qX)_{\ge1}$の部分加群である．
前補題とnoetherian性からこれは有限生成で，有限次元の次数0部分を補えば
$N^j(X)/gN^j(X)$全体も有限生成となる．前補題(1)より$N^j(X)$は有限生成である．
この議論では次数0の列の左端を誤って零と置かない．正次数の単射と有限な低次数の補正で十分である．

$F_qX$のcohomologyの範囲を$[a,b]$とすると，$C^j(F_qX)=0$は
全$n\ge0$について$j\notin[a,b+1]$で成立する．
また$\RHom(e,X)$も有界である．両者の範囲の外では
\eqref{eq:weak-central-exact}の$C^j$が零なので，$g:N^j_{n-1}\to N^j_n$は全射である．
$N^j_0=0$を初期値として，その範囲外の全成分が零になる．
ゆえに非零の$j$は有限個である．

\emph{関手の構成と核．}
正のdg orbit代数$\mathcal R_q=\bigoplus_{n\ge0}\RHom(e,s^ne)$を取る．
第一段階により，これはcohomological degree 0にしかcohomologyを持たない．
good truncationを介するdg代数のquasi-isomorphism
$\mathcal R_q\leftarrow\tau_{\le0}\mathcal R_q\to R_q$でmoduleを運ぶと，
$\bigoplus_{n\ge0}\RHom(e,s^nX)$が自然な$R_q$-moduleの導来対象を定め，
そのcohomologyは$N^j(X)$である．上の有限生成性と有界性から
$\Db(\gr R_q)$に入り，exactな商関手$\pi$を施して$\Phi_q$を得る．
有限個の低次数を除く内部shiftの同定が$\Phi_qs\simeq(1)\Phi_q$を与える．
cohomology公式と核の主張は命題\ref{prop:weak-coherent-kernel}および
定理\ref{thm:weak-orbit-kernel}による．genuine性，orbitによる全圏の生成，
標的の有限大域次元は，この構成には使っていない．
\end{proof}

\begin{remark}[三つの有限性の区別]\label{rem:three-finiteness}
この定理は，反標準\emph{環}のnoetherian性，個々の対象の次数付きExtの\emph{有限生成性}，
そのcohomological degreeの\emph{有界性}を順に証明している．
$g$可逆な非零加群には$N/gN=0$となる例があるが，ここで使う$N^j(X)$は
内部次数が非負の切断加群であり，下に有界性がその誤用を排除する．
また$\ker\Phi_q$の同定だけでは，誘導関手の完全忠実性・本質的全射性は従わない．
\end{remark}

\subsection{古典的特異モデル上のgluingとしての商}
$\bar i:E\hookrightarrow\overline S$を反標準収縮後の境界とし，
\[
 \overline{\cB}=\OO_{\overline S}^{\perp}\subset\Db(\coh\overline S),\qquad
 \mathcal M_q(G)=\RHom_E(L\bar i^*G,q)
\]
と置く．境界は特異点を通らず，Cartier divisorなので$L\bar i^*$は有界coherent複体を保つ．
$\overline{\cD}_q$を，$\overline{\cB}$にexceptional対象$\bar e$を
$\RHom(\bar e,G)=0$，$\RHom(G,\bar e)=\mathcal M_q(G)$で付けた
pretriangulated dg gluingの冪等完備化と定義する．
これは定義であり，ある非可換abel圏の有界導来圏であることを先取りしない．

\begin{theorem}[収縮圏の特異gluing表示]\label{thm:singular-gluing-target}
標準弱楕円gluingのdg商は
\[
 \boxed{(\cD_q/\mathcal N_q)^\natural\simeq\overline{\cD}_q.}
\]
この同値は$e$を$\bar e$へ送り，共通成分では$Rf_*:\cB\to\overline{\cB}$から誘導される．
$q=\OO_E$では$\overline{\cD}_{\OO_E}\simeq\Db(\coh\overline S)$である．
\end{theorem}
\begin{proof}
反標準収縮$f:S\to\overline S$は射影的で，幾何学的fiberの次元は高々1，
Du Val特異点の有理性から$Rf_*\OO_S=\OO_{\overline S}$である．
従ってBondal--Kapranov--Schechtman\cite[Theorem 2.14, \S2F]{BKSWalls}の
\emph{既知の}局所化定理が適用され，
\[
 \Db(\coh S)/\mathcal N_f\simeq\Db(\coh\overline S)
\]
をdg圏として得る．この定理は任意次元fiberについての一般局所化予想を必要としない．

$Rf_*$は構造対象を保存し，
$\RHom(\OO_S,G)=\RHom(\OO_{\overline S},Rf_*G)$である．
したがって$\cB/\mathcal N_f\simeq\overline{\cB}$となる．
具体的には，$\overline{\cB}$の対象を局所化定理によって持ち上げれば，大域切断の消滅から
その持ち上げも$\cB$に入る．$\cB$の二対象間の局所化のroofでは，coneが
$\mathcal N_f\subset\cB$に入るため，中間対象も$\cB$に取れる．これが完全忠実性を与える．

$f$は$E$の近傍で同型なので
$Li^*G\simeq L\bar i^*Rf_*G$であり，gluing moduleは
$\cB/\mathcal N_f$上で$\mathcal M_q$に一致する．
命題\ref{prop:zero-dg-quotient}のgluingとdg商の両立性を使えば所望の同定を得る．

$q=\OO_E$のとき，$\overline S$は射影Gorenstein有理曲面で
$\omega_{\overline S}=\OO_{\overline S}(-E)$である．
$G\in\overline{\cB}$についてGrothendieck双対性から
$\RHom(G,\omega_{\overline S})=R\Gamma(G)^*[-2]=0$であり，
Cartier divisorのtriangleが
$\RHom(G,\OO_{\overline S})\simeq\RHom_E(L\bar i^*G,\OO_E)$を与える．
また$\OO_{\overline S}$はexceptionalで，大域切断の評価triangleにより
$\Db(\coh\overline S)$は$\overline{\cB}$とこの対象で生成される．
従って上記gluingはその有界coherent導来圏に一致する．
\end{proof}

\begin{corollary}[零根壁の商圏には高次自己Extが残る]\label{cor:singular-gluing-unbounded}
$S$がgenuineでないなら，任意の$q$について$\mathcal Q_q$には
$\Hom(W,W[j])\ne0$となる$j>0$が無限にある対象$W$が存在する．
特に$\mathcal Q_q$はproper dg圏ではなく，滑らかなsource $\cD_q$と抽象的にも導来同値でない．
\end{corollary}
\begin{proof}
$\overline S$の特異点$x$を選び，そのideal sheafを$I_x$とする．
$R\Gamma(\OO_{\overline S})=k=R\Gamma(k(x))$なので$R\Gamma(I_x)=0$で，
$I_x\in\overline{\cB}$である．$\overline S$はGorenstein曲面なので
$\Ext^j(k(x),\OO_{\overline S})=0$ $(j>2)$．
$0\to I_x\to\OO_{\overline S}\to k(x)\to0$をRHomの両変数へ用いると
\[
 \Ext^j(I_x,I_x)\simeq\Ext^j(k(x),k(x))\quad(j>2)
\]
を得る．非正則局所環での剰余体の最小自由解消は無限に続き，
右辺は無限個の次数で非零である．支持が一点なので局所Extと大域Extは一致する．
$\overline{\cB}$は全ての$\overline{\cD}_q$へ完全忠実に入るから，
その中の$I_x$の像を$W$とすればよい．
source $\cD_q$では任意の二対象間のRHomの全cohomologyが有限次元であり，
この性質は三角圏同値で保存されるので最後の主張を得る．
\end{proof}
この系の標的は\emph{構成した}特異gluing $\overline{\cD}_q$である．
$q\ne\OO_E$でこれを$\Db(\qgr R_q)$と同定したことにはならない．

\subsection{残る比較問題の正確な形}\label{subsec:coherent-future}
本節によって，標準弱楕円gluingのcoherent有限生成性・有界性，
自然関手の全核，およびdg商の具体的な特異gluing表示が得られた．
一般の$q$について残る核心は，そこから誘導される自然関手
\[
 \overline\Phi_q:\overline{\cD}_q\longrightarrow\Db(\qgr R_q)
\]
の完全忠実性と本質的全射性である．
$q=\OO_E$では，反標準環の$qgr$が$\coh\overline S$であることと射影公式により，
$\Phi_q$は$Rf_*$に一致するので，局所化同値が成立する．
反標準orbitで生成される完全部分圏上では同値を証明できる
（定理\ref{thm:orbit-core-equivalence}）．しかし零指数壁では
この部分圏を除いた残余商$\mathfrak S_q$が非零である
（定理\ref{thm:residual-singularity}）．
従って非自明$q$の残りの比較は，特異点に対応する
非perfect対象とorbitの間のgluingおよび全高次Extを制御する問題である．
一つの$A_1$壁から始め，特異モデル側のideal sheafの像を
反標準$qgr$側で計算する．

指定Rains heartでの収縮対象の位置，壁上のAS模型との同定，局所・完成代数，
曲面族に対する底変更と降下は別の課題として残る．
本節の有限生成の補題，古典的な局所化定理，dg gluingの一般論自体の新規性は主張しない．
追加的な内容は，それらを同じ$(S,E,q,e,s)$に適用し，零指数の全核，
coherent反標準関手，特異gluingの三者を自然関手つきで接続した点にある．
